\subsection{Optimizer-Induced Neural Collapse and Feature Geometry}
\label{subsec:optimizer-geometry}

Neural Collapse provides a geometric description of the terminal phase of
supervised training, in which penultimate features, class means, and classifier
weights approach a highly structured configuration
\citep{papyan2020prevalence}. This perspective is useful for reliability because
it separates accuracy from representation geometry: two classifiers can achieve
similar ID performance while differing in within-class compactness, class-mean
separation, classifier--feature alignment, feature-norm structure, and covariance
shape. Feature-based OOD detectors operate precisely on these geometric
quantities, so changes in Neural Collapse or related feature geometry can affect
detector behavior even when classification accuracy is comparable.

Recent work shows that this geometry is not optimizer invariant.
\citet{zhao2026optimizer} demonstrate that optimizer choice and weight-decay
coupling affect the emergence of Neural Collapse, with SGD, Adam, and AdamW
exhibiting different collapse behavior. This is closely related to the distinction
between coupled and decoupled weight decay \citep{loshchilov2019decoupled}. In
SGD-like optimizers, coupled and decoupled decay can be relatively similar in
their practical effect, whereas in adaptive optimizers the distinction can change
the implicit bias of training. Our work builds on this optimizer-to-geometry
view, but studies the next step: how geometry differences induced by SGD, SGDW,
Adam, and AdamW are inherited by downstream OOD detector families.

\subsection{Feature-Based OOD Detectors as Geometry Readouts}
\label{subsec:feature-based-ood-readouts}

OOD detection methods assign scores intended to separate ID-like inputs from
unreliable or out-of-distribution inputs. Logit-based scores such as MSP,
MaxLogit, and Energy read output confidence, logit scale, and margin structure
\citep{hendrycks2017baseline,liu2020energy}. They are useful and often strong
baselines, but they do not directly model the geometry of the penultimate feature
space.

Feature-based detectors instead read different geometric channels of the learned
representation. Mahalanobis scoring fits class means and covariance estimates and
uses a precision-weighted distance to the closest class center
\citep{lee2018simple}. kNN scoring compares a query feature to a bank of ID
training features and is therefore sensitive to local neighborhood spacing,
feature normalization, and local density structure \citep{sun2022knn}. DDU-style
Gaussian density scores fit class-conditional feature densities and read both
distance-to-mean and covariance-volume information \citep{mukhoti2021ddu}.
Angular detectors such as CTM use cosine similarity to class-typical feature
directions \citep{nguyen2023ctm}, while Neural-Collapse-based detectors such as
NeCo use prototype, class-mean, and principal-component structure
\citep{benammar2023neco}. These detectors are therefore not interchangeable
scoring rules: each one reads a different projection of the optimizer-induced
feature geometry.

This readout view is central to our paper. If an optimizer changes feature norms,
raw distance detectors can change even when angular structure is preserved. If it
changes covariance anisotropy or conditioning, Gaussian density and
Mahalanobis-style scores can change even when class means remain separated. If it
changes class-mean directions, classifier alignment, or residual subspaces,
angular and collapse-based scores can change even when raw distance-based scores
appear stable. We therefore organize detectors by the geometry channel they read,
rather than by treating all OOD scores as a single benchmark list.

\subsection{Normalization, Covariance, and Subspace Controls}
\label{subsec:normalization-covariance-controls}

Feature-space preprocessing can substantially change the behavior of
feature-based detectors. In particular, L2 normalization removes radial
feature-norm variation and turns Euclidean comparisons into angular comparisons.
This is important because normalized kNN is often a strong practical variant, and
normalization can also improve Mahalanobis-style scoring when raw feature norms
distort covariance estimates. In our framework, however, the purpose of raw and
normalized variants is not only to compare detector performance. The comparison
itself is diagnostic: raw Mahalanobis and raw kNN expose radial or norm-coupled
geometry, while L2-normalized versions test whether the detector gap remains once
that radial channel is removed.

Covariance-side controls play a complementary role. Mahalanobis and
DDU-style Gaussian density scores depend on covariance estimates, and these
estimates can be sensitive to feature dimension, sample size, anisotropy, and
conditioning. Tied covariance, diagonal covariance, and shrinkage covariance
therefore provide different probes of the same feature distribution. If shrinkage
or diagonalization recovers a detector gap, the failure may be tied to covariance
estimation or ill-conditioning. If the gap remains, the mismatch is more likely
to involve broader representation geometry, such as class overlap, radial
variation, or subspace structure.

PCA-based controls provide another way to separate geometry channels. Projecting
features onto a low-dimensional principal subspace can remove nuisance
directions, while residual-space scores can test whether OOD samples occupy
directions not used by the ID class structure. These transformations are not
introduced here as new detector proposals. Instead, we use them as diagnostic
interventions that help localize whether optimizer-induced detector changes are
carried by radial, covariance, angular, or subspace channels.

\subsection{Representation-Centric Detector Selection}
\label{subsec:representation-centric-detector-selection}

Recent representation-centric analyses argue that detector rankings are not fixed
properties of score formulas. They can change with architecture, training
paradigm, source dataset, representation geometry, and OOD severity
\citep{olivares2025systematic}. This perspective is important for our setting
because optimizer choice is a training-side intervention that can reshape the
representation seen by every downstream detector. A detector that is competitive
under one optimizer-induced geometry may become less reliable under another, not
because the detector formula is intrinsically weak, but because the geometry it
expects has been distorted, masked, or replaced by a different feature regime.

Our work follows this representation-centric view but focuses on a different
source of variation. Rather than varying detector design or training paradigm
broadly, we isolate optimizer choice and weight-decay coupling as mechanisms that
shape penultimate geometry. We then ask whether Neural Collapse metrics,
feature-norm statistics, within-class dispersion, class separation, covariance
spectra, effective rank, and local feature-bank structure can explain or predict
which detector family will work. This turns detector selection into a
geometry-aware diagnostic problem: before choosing a feature-based OOD detector,
one should first inspect the geometry that the trained model exposes to that
detector.

\subsection{Positioning of Our Work}
\label{subsec:positioning}

Prior work provides the components of our argument. Neural Collapse studies show
that final-layer and penultimate features can organize into measurable geometric
structures. Optimizer studies show that this geometry can depend on the training
rule and the coupling of weight decay. Feature-based OOD detectors rely on class
means, covariance, neighborhoods, densities, prototypes, or subspaces.
Representation-centric benchmarks show that detector rankings shift when the
learned representation changes.

The missing link is the optimizer--geometry--detector pathway. We study this
pathway directly by treating SGD, SGDW, Adam, and AdamW as geometry interventions
and by evaluating how their induced penultimate representations are read by
logit-based, distance-based, density-based, neighborhood-based, angular, and
Neural-Collapse-based detectors. The contribution is not another OOD scoring rule
and not a generic optimizer leaderboard. It is a geometry--detector compatibility
framework: optimizer choice should be evaluated jointly with the penultimate
geometry it induces and the detector family that will read that geometry.
